%==========================================================
%  CHAPITRE 3 — ÉTUDE DES FONCTIONS
%  Graphes intégrés avec pgfplots + tikz
%==========================================================

\chapter{\'{E}tude des Fonctions}

\begin{rappel}[M\'{e}thode g\'{e}n\'{e}rale d'\'{e}tude d'une fonction]
  L'\'{e}tude compl\`{e}te d'une fonction $f$ comprend :
  \begin{enumerate}
    \item Domaine de d\'{e}finition $D_f$
    \item Parit\'{e}, p\'{e}riodicit\'{e}, sym\'{e}tries
    \item Limites aux bornes et asymptotes
    \item D\'{e}riv\'{e}e $f'$, signe et tableau de variations
    \item Convexit\'{e} : d\'{e}riv\'{e}e seconde $f''$ et points d'inflexion
    \item Trac\'{e} de la courbe repr\'{e}sentative $\mathcal{C}_f$
  \end{enumerate}
\end{rappel}

%----------------------------------------------------------
\section{Sym\'{e}tries d'une courbe}
%----------------------------------------------------------

\subsection{Axe de sym\'{e}trie}

\begin{definition}[Axe de sym\'{e}trie]
  La droite d'\'{e}quation $x = a$ est un \textbf{axe de sym\'{e}trie}
  de la courbe $\mathcal{C}_f$ si et seulement si :
  \[
    \forall x \in D_f \;:\quad (2a - x) \in D_f
    \quad \text{et} \quad
    f(2a - x) = f(x)
  \]
\end{definition}

\begin{remarque}
  Cas particulier : si $a = 0$, la condition devient $f(-x) = f(x)$,
  ce qui signifie que $f$ est une \textbf{fonction paire}
  (sym\'{e}trique par rapport \`{a} l'axe des ordonn\'{e}es).
\end{remarque}

\begin{exemple}[R\'{e}solu : axe de sym\'{e}trie]
  $f(x) = x^2 - 4x + 5$.
  On calcule $f(4 - x) = (4-x)^2 - 4(4-x) + 5 = x^2 - 4x + 5 = f(x)$.

  Donc la droite $x = 2$ est un \textbf{axe de sym\'{e}trie} de $\mathcal{C}_f$.

  \medskip
  \begin{center}
  \begin{tikzpicture}
    \begin{axis}[
      axis lines  = center,
      xlabel      = {$x$},
      ylabel      = {$y$},
      xmin = -0.5, xmax = 4.5,
      ymin = -0.5, ymax = 7,
      xtick       = {0,1,2,3,4},
      ytick       = {0,1,2,3,4,5,6},
      grid        = major,
      grid style  = {gray!20},
      width = 9cm, height = 6cm,
    ]
      % Courbe f(x) = x²-4x+5
      \addplot[navy, very thick, domain=-0.3:4.3, samples=120]
        {x^2 - 4*x + 5};
      \addlegendentry{$f(x)=x^2-4x+5$}

      % Axe de symétrie x=2
      \addplot[burgundy, dashed, thick]
        coordinates {(2,-0.5)(2,7)};
      \addlegendentry{Axe $x=2$}

      % Sommet
      \addplot[mark=*, mark size=3pt, burgundy]
        coordinates {(2,1)};
      \node[right, color=burgundy, font=\small] at (axis cs:2,1) {$(2,\,1)$};
    \end{axis}
  \end{tikzpicture}
  \end{center}
\end{exemple}

%----------------------------------------------------------
\subsection{Centre de sym\'{e}trie}
%----------------------------------------------------------

\begin{definition}[Centre de sym\'{e}trie]
  Le point $I(a,\,b)$ est un \textbf{centre de sym\'{e}trie}
  de la courbe $\mathcal{C}_f$ si et seulement si :
  \[
    \forall x \in D_f \;:\quad (2a - x) \in D_f
    \quad \text{et} \quad
    f(2a - x) + f(x) = 2b
  \]
\end{definition}

\begin{remarque}
  Cas particulier : si $I = (0, 0)$, la condition devient $f(-x) = -f(x)$
  — $f$ est une \textbf{fonction impaire} (sym\'{e}trique par rapport \`{a} l'origine).
\end{remarque}

\begin{methode}[V\'{e}rifier qu'un point est centre de sym\'{e}trie]
  \begin{enumerate}
    \item Calculer directement $f(2a - x) + f(x)$.
    \item Si le r\'{e}sultat vaut $2b$ (constante), alors $I(a,b)$
          est centre de sym\'{e}trie.
  \end{enumerate}
\end{methode}

\begin{exemple}[R\'{e}solu : centre de sym\'{e}trie]
  $f(x) = x^3 - 3x$ sur $\R$.
  V\'{e}rifions que $I(0,\,0)$ est centre de sym\'{e}trie.

  \medskip
  $f(-x) = -x^3 + 3x = -f(x)$,
  donc $f(-x) + f(x) = 0 = 2\times 0$ : $I(0,0)$ est bien centre de sym\'{e}trie.

  \medskip
  \begin{center}
  \begin{tikzpicture}
    \begin{axis}[
      axis lines  = center,
      xlabel      = {$x$},
      ylabel      = {$y$},
      xmin = -2.5, xmax = 2.5,
      ymin = -4,   ymax = 4,
      grid        = major,
      grid style  = {gray!20},
      width = 9cm, height = 7cm,
    ]
      % Courbe f(x) = x³ - 3x
      \addplot[navy, very thick, domain=-2.2:2.2, samples=150]
        {x^3 - 3*x};
      \addlegendentry{$f(x)=x^3-3x$}

      % Centre de symétrie O(0,0)
      \addplot[mark=*, mark size=4pt, burgundy]
        coordinates {(0,0)};
      \node[above right, color=burgundy, font=\small\bfseries]
        at (axis cs:0,0) {$I(0,0)$};

      % Deux points symétriques
      \addplot[mark=square*, mark size=3pt, forest]
        coordinates {(1,-2)(-1,2)};
      \node[right, color=forest, font=\small] at (axis cs:1,-2)  {$(1,-2)$};
      \node[left,  color=forest, font=\small] at (axis cs:-1, 2) {$(-1,2)$};

      % Flèche de symétrie
      \draw[->, forest, dashed] (axis cs:1,-2) -- (axis cs:-1,2);
    \end{axis}
  \end{tikzpicture}
  \end{center}
\end{exemple}

%----------------------------------------------------------
\section{Convexit\'{e} et point d'inflexion}
%----------------------------------------------------------

\begin{definition}[Convexit\'{e} et concavit\'{e}]
  Soit $f$ deux fois d\'{e}rivable sur un intervalle $I$.
  \begin{itemize}
    \item $\mathcal{C}_f$ est \textbf{convexe} sur $I$ si elle est situ\'{e}e
          \textbf{au-dessus} de chacune de ses tangentes :
          \( \forall x \in I,\quad f''(x) \geq 0 \)
    \item $\mathcal{C}_f$ est \textbf{concave} sur $I$ si elle est situ\'{e}e
          \textbf{au-dessous} de chacune de ses tangentes :
          \( \forall x \in I,\quad f''(x) \leq 0 \)
  \end{itemize}
\end{definition}

\begin{definition}[Point d'inflexion]
  Un \textbf{point d'inflexion} est un point o\`{u} $\mathcal{C}_f$
  change de concavit\'{e}.

  \medskip
  \begin{center}
  \renewcommand{\arraystretch}{1.8}
  \begin{tabular}{|p{6.5cm}|p{6.5cm}|}
    \hline
    \textbf{\color{navy}Condition}
      & \textbf{\color{navy}Conclusion} \\
    \hline\hline
    $f''(x_0) = 0$ et $f''$ \textbf{change de signe} en $x_0$
      & $\mathcal{C}_f$ admet un \textbf{point d'inflexion} en $x_0$ \\
    \hline
    $f''(x_0) = 0$ mais \textbf{sans} changement de signe
      & Pas de point d'inflexion en $x_0$ \\
    \hline
  \end{tabular}
  \end{center}
\end{definition}

\begin{attention}
  $f''(x_0) = 0$ \textbf{ne suffit pas} !
  Il faut \textbf{obligatoirement} v\'{e}rifier le changement de signe de $f''$.
\end{attention}

\begin{methode}[\'{E}tudier la convexit\'{e} et trouver les points d'inflexion]
  \begin{enumerate}
    \item Calculer $f''(x)$.
    \item \'{E}tudier le signe de $f''(x)$ sur chaque intervalle.
    \item Les z\'{e}ros o\`{u} $f''$ change de signe $\Rightarrow$ points d'inflexion.
    \item Calculer les coordonn\'{e}es : $\bigl(x_0,\, f(x_0)\bigr)$.
  \end{enumerate}
\end{methode}

\begin{exemple}[R\'{e}solu : convexit\'{e} et point d'inflexion]
  $f(x) = x^3 - 3x^2 + 2$ sur $\R$.

  \medskip
  $f'(x) = 3x^2 - 6x$ \quad et \quad $f''(x) = 6x - 6 = 6(x-1)$.

  \medskip
  \begin{center}
  \renewcommand{\arraystretch}{1.6}
  \begin{tabular}{|c|ccccc|}
    \hline
    $x$           & $-\infty$ && $1$ && $+\infty$ \\
    \hline
    $f''(x)$      && $-$ & $0$ & $+$ & \\
    \hline
    Convexit\'{e} & \multicolumn{2}{c}{Concave $\frown$}
                  && \multicolumn{2}{c|}{Convexe $\smile$} \\
    \hline
  \end{tabular}
  \end{center}

  \medskip
  $f''$ s'annule et change de signe en $x_0 = 1$.
  Le \textbf{point d'inflexion} est $(1,\, f(1)) = (1,\, 0)$.

  \medskip
  \begin{center}
  \begin{tikzpicture}
    \begin{axis}[
      axis lines  = center,
      xlabel      = {$x$},
      ylabel      = {$y$},
      xmin = -0.8, xmax = 3.2,
      ymin = -3,   ymax = 4,
      xtick       = {-0,1,2,3},
      grid        = major,
      grid style  = {gray!20},
      width = 10cm, height = 7cm,
    ]
      % Courbe f(x) = x³ - 3x² + 2
      \addplot[navy, very thick, domain=-0.6:3.1, samples=150]
        {x^3 - 3*x^2 + 2};
      \addlegendentry{$f(x)=x^3-3x^2+2$}

      % Tangente au point d'inflexion (1,0) : f'(1)=-3 → y=-3(x-1)
      \addplot[burgundy, dashed, thick, domain=-0.2:2.5]
        {-3*(x-1)};
      \addlegendentry{Tangente au pt d'inflexion}

      % Point d'inflexion
      \addplot[mark=square*, mark size=4pt, burgundy]
        coordinates {(1,0)};
      \node[above right, color=burgundy, font=\small\bfseries]
        at (axis cs:1,0) {Inflexion $(1,0)$};

      % Annotation concave/convexe
      \node[color=navy!70, font=\small\bfseries] at (axis cs:0.3,2.5)
        {Concave $\frown$};
      \node[color=forest, font=\small\bfseries] at (axis cs:2.5,-1.5)
        {Convexe $\smile$};
    \end{axis}
  \end{tikzpicture}
  \end{center}
\end{exemple}
\newpage
\
%----------------------------------------------------------
\section{Branches infinies et asymptotes}
%----------------------------------------------------------

\begin{definition}[Types d'asymptotes]
  \begin{center}
  \renewcommand{\arraystretch}{2}
  \begin{tabular}{|p{5.5cm}|p{7cm}|}
    \hline
    \textbf{\color{navy}Condition}
      & \textbf{\color{navy}Conclusion} \\
    \hline\hline
    $\displaystyle\lim_{x \to a} f(x) = \pm\infty$
      & Asymptote \textbf{verticale} $x = a$ \\
    \hline
    $\displaystyle\lim_{x \to \pm\infty} f(x) = b \in \R$
      & Asymptote \textbf{horizontale} $y = b$ \\
    \hline
    $\displaystyle\lim_{x \to \pm\infty} \frac{f(x)}{x} = a \neq 0$
    et $\lim [f(x)-ax] = b \in \R$
      & Asymptote \textbf{oblique} $y = ax+b$ \\
    \hline
  \end{tabular}
  \end{center}
\end{definition}

\begin{methode}[Chercher une asymptote oblique en $\pm\infty$]
  \begin{enumerate}
    \item Calculer $a = \displaystyle\lim_{x \to \pm\infty} \dfrac{f(x)}{x}$.
    \item Si $a \in \R^*$, calculer $b = \displaystyle\lim_{x \to \pm\infty}[f(x) - ax]$.
    \item Si $b \in \R$ $\Rightarrow$ asymptote oblique $y = ax+b$.
    \item Si $b = \pm\infty$ $\Rightarrow$ branche parabolique de direction $y=ax$.
  \end{enumerate}
\end{methode}

\begin{exemple}[R\'{e}solu : asymptotes verticale et oblique]
  $f(x) = \dfrac{x^2}{x-1}$ sur $\R\setminus\{1\}$.

  \medskip
  \textbf{A.V. :} $\displaystyle\lim_{x\to 1} f(x) = \pm\infty$
  $\Rightarrow$ asymptote verticale $x = 1$.

  \medskip
  \textbf{A.O. :} $a = \displaystyle\lim_{x\to+\infty}\dfrac{f(x)}{x}
  = \lim\dfrac{x}{x-1} = 1$. \quad
  $b = \lim[f(x)-x] = \lim\dfrac{x^2 - x(x-1)}{x-1} = \lim\dfrac{x}{x-1} = 1$.

  Donc $y = x + 1$ est asymptote oblique.

  \medskip
  \begin{center}
  \begin{tikzpicture}
    \begin{axis}[
      axis lines  = center,
      xlabel      = {$x$},
      ylabel      = {$y$},
      xmin = -3, xmax = 5,
      ymin = -6, ymax = 9,
      grid        = major,
      grid style  = {gray!20},
      width = 10cm, height = 8cm,
      restrict y to domain = -8:10,
    ]
      % Branche gauche
      \addplot[navy, very thick, domain=-2.8:0.75, samples=200]
        {x^2/(x-1)};

      % Branche droite
      \addplot[navy, very thick, domain=1.3:4.8, samples=200]
        {x^2/(x-1)};
      \addlegendentry{$f(x)=\dfrac{x^2}{x-1}$}

      % Asymptote verticale x=1
      \addplot[burgundy, dashed, thick]
        coordinates {(1,-6)(1,9)};
      \addlegendentry{A.V. : $x=1$}

      % Asymptote oblique y = x+1
      \addplot[forest, dashed, thick, domain=-3:5]
        {x+1};
      \addlegendentry{A.O. : $y=x+1$}
    \end{axis}
  \end{tikzpicture}
  \end{center}
\end{exemple}

%----------------------------------------------------------
\section{Fonction r\'{e}ciproque}
%----------------------------------------------------------

\begin{definition}[Fonction r\'{e}ciproque]
  Si $f$ est \textbf{continue et strictement monotone} sur $I$,
  alors $f$ admet une \textbf{fonction r\'{e}ciproque} $f^{-1}$
  d\'{e}finie sur $f(I)$, v\'{e}rifiant :
  \[
    f^{-1}(x) = y \;\iff\; f(y) = x
  \]
\end{definition}

\begin{propriete}[Propri\'{e}t\'{e}s de $f^{-1}$]
  \begin{itemize}
    \item $f^{-1}$ est \textbf{continue} sur $f(I)$.
    \item $f^{-1}$ varie dans le \textbf{m\^{e}me sens} que $f$.
    \item Si $f'$ ne s'annule pas sur $I$ :
    \[
      \bigl(f^{-1}\bigr)'(x) = \frac{1}{f'\bigl(f^{-1}(x)\bigr)}
      \qquad \text{et} \qquad
      \bigl(f^{-1}\bigr)'(y_0) = \frac{1}{f'(x_0)}
      \quad \text{avec } y_0 = f(x_0)
    \]
  \end{itemize}
\end{propriete}

\begin{methode}[Trouver l'expression de $f^{-1}(x)$]
  \begin{enumerate}
    \item Poser $f(y) = x$.
    \item \textbf{R\'{e}soudre} cette \'{e}quation en $y$.
    \item La solution $y$ est $f^{-1}(x)$.
  \end{enumerate}
\end{methode}

\begin{exemple}[R\'{e}solu : $f$ et $f^{-1}$ sym\'{e}triques par rapport \`{a} $y=x$]
  $f(x) = e^x$ sur $\R$, donc $f^{-1}(x) = \ln x$ sur $]0,+\infty[$.

  \medskip
  \begin{center}
  \begin{tikzpicture}
    \begin{axis}[
      axis lines  = center,
      xlabel      = {$x$},
      ylabel      = {$y$},
      xmin = -2.5, xmax = 4,
      ymin = -2.5, ymax = 4,
      grid        = major,
      grid style  = {gray!20},
      width = 9cm, height = 9cm,
    ]
      % f(x) = e^x
      \addplot[navy, very thick, domain=-2.4:1.38, samples=120]
        {exp(x)};
      \addlegendentry{$f(x)=e^x$}

      % f⁻¹(x) = ln(x)
      \addplot[forest, very thick, domain=0.09:4, samples=120]
        {ln(x)};
      \addlegendentry{$f^{-1}(x)=\ln x$}

      % Première bissectrice y = x
      \addplot[burgundy, dashed, thick, domain=-2:4]
        {x};
      \addlegendentry{Bissectrice $y=x$}

      % Point commun (0,1) et (1,0)
      \addplot[mark=*, mark size=3pt, navy]
        coordinates {(0,1)};
      \node[left, font=\small, navy] at (axis cs:0,1) {$(0,1)$};

      \addplot[mark=*, mark size=3pt, forest]
        coordinates {(1,0)};
      \node[below, font=\small, forest] at (axis cs:1,0) {$(1,0)$};
    \end{axis}
  \end{tikzpicture}
  \end{center}
\end{exemple}

\begin{propriete}[Correspondance des asymptotes entre $f$ et $f^{-1}$]
  \begin{center}
  \renewcommand{\arraystretch}{1.8}
  \begin{tabular}{|p{5.5cm}|p{5.5cm}|}
    \hline
    \textbf{\color{navy}$\mathcal{C}_f$ admet\ldots}
      & \textbf{\color{navy}$\mathcal{C}_{f^{-1}}$ admet\ldots} \\
    \hline\hline
    A. verticale $x = a$     & A. horizontale $y = a$ \\
    \hline
    A. horizontale $y = a$   & A. verticale $x = a$ \\
    \hline
    A. oblique $y = ax + b$  & A. oblique $y = \dfrac{1}{a}x - \dfrac{b}{a}$ \\
    \hline
    Tangente verticale       & Tangente horizontale \\
    \hline
    Tangente horizontale     & Tangente verticale \\
    \hline
  \end{tabular}
  \end{center}
\end{propriete}

\newpage
%----------------------------------------------------------
\section*{Exercices du chapitre}
\addcontentsline{toc}{section}{Exercices : \'{E}tude des fonctions}
%----------------------------------------------------------

\begin{exercice}[1 : Sym\'{e}tries]
  Pour chaque fonction, d\'{e}terminer si la courbe poss\`{e}de
  un axe ou un centre de sym\'{e}trie :
  \begin{enumerate}
    \item $f(x) = x^4 - 2x^2 + 1$
    \item $f(x) = x^3 - 3x$
    \item $f(x) = \dfrac{1}{x-2} + 3$
  \end{enumerate}
\end{exercice}

\begin{exercice}[2 : Convexit\'{e} et points d'inflexion]
  \'{E}tudier la convexit\'{e} et d\'{e}terminer les \'{e}ventuels points d'inflexion :
  \begin{enumerate}
    \item $f(x) = x^3 - 6x^2 + 9x$ sur $\R$
    \item $f(x) = e^x - x$ sur $\R$
    \item $f(x) = \ln x$ sur $]0,+\infty[$
  \end{enumerate}
\end{exercice}

\begin{exercice}[3 : Asymptotes]
  \'{E}tudier les branches infinies de :
  \begin{enumerate}
    \item $f(x) = \dfrac{x^2-1}{x+1}$ \quad (asymptote oblique ?)
    \item $f(x) = \dfrac{2x}{x-3}$ \quad (asymptotes verticale et horizontale)
    \item $f(x) = x + \sqrt{x^2+1}$ \quad (asymptote oblique en $\pm\infty$)
  \end{enumerate}
\end{exercice}

\begin{exercice}[4 : Fonction r\'{e}ciproque]
  D\'{e}terminer $f^{-1}$ et son domaine :
  \begin{enumerate}
    \item $f(x) = 3x - 5$ sur $\R$
    \item $f(x) = x^2 + 1$ sur $[0,+\infty[$
    \item $f(x) = e^{x+1}$ sur $\R$
    \item $f(x) = \dfrac{x+1}{x-1}$ sur $]1,+\infty[$
  \end{enumerate}
\end{exercice}

\begin{exercice}[5 : \'{E}tude compl\`{e}te avec trac\'{e}]
  Soit $f(x) = \dfrac{x^2}{x-1}$ sur $\R\setminus\{1\}$.
  \begin{enumerate}
    \item Limites aux bornes et asymptotes.
    \item $f'(x)$, signe et tableau de variations.
    \item $f''(x)$, convexit\'{e} et point d'inflexion.
    \item Tracer $\mathcal{C}_f$ avec ses asymptotes.
  \end{enumerate}
\end{exercice}